% Options for packages loaded elsewhere
% Options for packages loaded elsewhere
\PassOptionsToPackage{unicode}{hyperref}
\PassOptionsToPackage{hyphens}{url}
\PassOptionsToPackage{dvipsnames,svgnames,x11names}{xcolor}
%
\documentclass[
  french,
  letterpaper,
  DIV=11,
  numbers=noendperiod]{scrreprt}
\usepackage{xcolor}
\usepackage{amsmath,amssymb}
\setcounter{secnumdepth}{5}
\usepackage{iftex}
\ifPDFTeX
  \usepackage[T1]{fontenc}
  \usepackage[utf8]{inputenc}
  \usepackage{textcomp} % provide euro and other symbols
\else % if luatex or xetex
  \usepackage{unicode-math} % this also loads fontspec
  \defaultfontfeatures{Scale=MatchLowercase}
  \defaultfontfeatures[\rmfamily]{Ligatures=TeX,Scale=1}
\fi
\usepackage{lmodern}
\ifPDFTeX\else
  % xetex/luatex font selection
\fi
% Use upquote if available, for straight quotes in verbatim environments
\IfFileExists{upquote.sty}{\usepackage{upquote}}{}
\IfFileExists{microtype.sty}{% use microtype if available
  \usepackage[]{microtype}
  \UseMicrotypeSet[protrusion]{basicmath} % disable protrusion for tt fonts
}{}
\makeatletter
\@ifundefined{KOMAClassName}{% if non-KOMA class
  \IfFileExists{parskip.sty}{%
    \usepackage{parskip}
  }{% else
    \setlength{\parindent}{0pt}
    \setlength{\parskip}{6pt plus 2pt minus 1pt}}
}{% if KOMA class
  \KOMAoptions{parskip=half}}
\makeatother
% Make \paragraph and \subparagraph free-standing
\makeatletter
\ifx\paragraph\undefined\else
  \let\oldparagraph\paragraph
  \renewcommand{\paragraph}{
    \@ifstar
      \xxxParagraphStar
      \xxxParagraphNoStar
  }
  \newcommand{\xxxParagraphStar}[1]{\oldparagraph*{#1}\mbox{}}
  \newcommand{\xxxParagraphNoStar}[1]{\oldparagraph{#1}\mbox{}}
\fi
\ifx\subparagraph\undefined\else
  \let\oldsubparagraph\subparagraph
  \renewcommand{\subparagraph}{
    \@ifstar
      \xxxSubParagraphStar
      \xxxSubParagraphNoStar
  }
  \newcommand{\xxxSubParagraphStar}[1]{\oldsubparagraph*{#1}\mbox{}}
  \newcommand{\xxxSubParagraphNoStar}[1]{\oldsubparagraph{#1}\mbox{}}
\fi
\makeatother


\usepackage{longtable,booktabs,array}
\newcounter{none} % for unnumbered tables
\usepackage{calc} % for calculating minipage widths
% Correct order of tables after \paragraph or \subparagraph
\usepackage{etoolbox}
\makeatletter
\patchcmd\longtable{\par}{\if@noskipsec\mbox{}\fi\par}{}{}
\makeatother
% Allow footnotes in longtable head/foot
\IfFileExists{footnotehyper.sty}{\usepackage{footnotehyper}}{\usepackage{footnote}}
\makesavenoteenv{longtable}
\usepackage{graphicx}
\makeatletter
\newsavebox\pandoc@box
\newcommand*\pandocbounded[1]{% scales image to fit in text height/width
  \sbox\pandoc@box{#1}%
  \Gscale@div\@tempa{\textheight}{\dimexpr\ht\pandoc@box+\dp\pandoc@box\relax}%
  \Gscale@div\@tempb{\linewidth}{\wd\pandoc@box}%
  \ifdim\@tempb\p@<\@tempa\p@\let\@tempa\@tempb\fi% select the smaller of both
  \ifdim\@tempa\p@<\p@\scalebox{\@tempa}{\usebox\pandoc@box}%
  \else\usebox{\pandoc@box}%
  \fi%
}
% Set default figure placement to htbp
\def\fps@figure{htbp}
\makeatother



\ifLuaTeX
\usepackage[bidi=basic,shorthands=off]{babel}
\else
\usepackage[bidi=default,shorthands=off]{babel}
\fi
\ifLuaTeX
  \usepackage{selnolig} % disable illegal ligatures
\fi


\setlength{\emergencystretch}{3em} % prevent overfull lines

\providecommand{\tightlist}{%
  \setlength{\itemsep}{0pt}\setlength{\parskip}{0pt}}



 


\KOMAoption{captions}{tableheading}
\makeatletter
\@ifpackageloaded{bookmark}{}{\usepackage{bookmark}}
\makeatother
\makeatletter
\@ifpackageloaded{caption}{}{\usepackage{caption}}
\AtBeginDocument{%
\ifdefined\contentsname
  \renewcommand*\contentsname{Table des matières}
\else
  \newcommand\contentsname{Table des matières}
\fi
\ifdefined\listfigurename
  \renewcommand*\listfigurename{Liste des Figures}
\else
  \newcommand\listfigurename{Liste des Figures}
\fi
\ifdefined\listtablename
  \renewcommand*\listtablename{Liste des Tables}
\else
  \newcommand\listtablename{Liste des Tables}
\fi
\ifdefined\figurename
  \renewcommand*\figurename{Figure}
\else
  \newcommand\figurename{Figure}
\fi
\ifdefined\tablename
  \renewcommand*\tablename{Table}
\else
  \newcommand\tablename{Table}
\fi
}
\@ifpackageloaded{float}{}{\usepackage{float}}
\floatstyle{ruled}
\@ifundefined{c@chapter}{\newfloat{codelisting}{h}{lop}}{\newfloat{codelisting}{h}{lop}[chapter]}
\floatname{codelisting}{Listing}
\newcommand*\listoflistings{\listof{codelisting}{Liste des Listings}}
\makeatother
\makeatletter
\makeatother
\makeatletter
\@ifpackageloaded{caption}{}{\usepackage{caption}}
\@ifpackageloaded{subcaption}{}{\usepackage{subcaption}}
\makeatother
\usepackage{bookmark}
\IfFileExists{xurl.sty}{\usepackage{xurl}}{} % add URL line breaks if available
\urlstyle{same}
\hypersetup{
  pdftitle={Pour une pédagogie perceptive},
  pdfauthor={Matthieu Gaudeau},
  pdflang={fr},
  colorlinks=true,
  linkcolor={blue},
  filecolor={Maroon},
  citecolor={Blue},
  urlcolor={Blue},
  pdfcreator={LaTeX via pandoc}}


\title{Pour une pédagogie perceptive}
\author{Matthieu Gaudeau}
\date{2026-01-01}
\begin{document}
\maketitle

\renewcommand*\contentsname{Table des matières}
{
\setcounter{tocdepth}{2}
\tableofcontents
}

\bookmarksetup{startatroot}

\chapter*{Preface}\label{preface}
\addcontentsline{toc}{chapter}{Preface}

\markboth{Preface}{Preface}

This is a Quarto book.

To learn more about Quarto books visit
\url{https://quarto.org/docs/books}.

\bookmarksetup{startatroot}

\chapter{}\label{section}

Ce livre est là parce qu'il m'a manqué. Il a manqué quand je me suis
formé à la Technique Alexander, puis quand j'ai commencé à la
transmettre à d'autres et que j'ai eu besoin de la parler.~

La Technique Alexander « fonctionnait », j'en faisais l'expérience, et
pourtant quelque chose ne se disait pas.

Ce quelque chose n'était pas un secret. Ce n'était pas non plus un
ineffable qui décourage l'enquête. C'était plutôt une lacune dans le
langage disponible : ce léger déplacement de monde qui accompagne
parfois la main d'un enseignant, cette façon dont une attention
différente rend le sol soudainement plus présent, le plafond plus
lointain, soi-même moins évident.

Ce livre est une tentative d'offrir ces mots. Des mots nés de la
nécessité de comprendre et de transmettre, et qui ne sont pas venus dans
la solitude de la lecture. Ils ont émergé dans des contextes d'échange,
de transmission, de projets où la Technique Alexander devait se dire à
d'autres disciplines, à d'autres langages.

Pour cela, j'ai dû aller chercher ailleurs. Du côté de la
phénoménologie, qui s'est donné pour tâche depuis Husserl de décrire
l'expérience avant qu'elle ne soit réduite à ses conditions causales. Du
côté de l'énaction et de la cognition incarnée, qui ont montré que
percevoir n'est pas enregistrer un monde extérieur mais le co-produire
dans l'acte même d'y être engagé. Du côté de la psychologie écologique
de Gibson, dont l'étrange concept d'\emph{affordance} --- ces
possibilités d'action que les choses offrent aux organismes --- éclaire
mieux que beaucoup d'autres l'expérience de quelqu'un qui retrouve
soudain son poids, sa verticalité, son sentiment d'espace.

Ces emprunts ne sont pas des confirmations théoriques. Ils répondent à
une nécessité : si la Technique Alexander transforme quelque chose dans
la façon dont on perçoit et dont on se meut, il faut une théorie de la
perception et du mouvement qui soit à la hauteur de cette
transformation. Les cadres disponibles dans la littérature alexandrienne
classique --- fondés sur des métaphores mécanistes ou des intuitions
évolutionnistes datées --- ne suffisent plus. Il nous faut une
phénoménologie du geste.

C'est ce que ce livre propose d'esquisser. Loin d'être une théorie
complète, il offre un cadre conceptuel, un vocabulaire, une façon de
regarder --- qui permette de parler de ce qui se passe quand quelqu'un
apprend la Technique Alexander, et peut-être, en retour, de mieux
comprendre ce que signifie percevoir, se mouvoir, faire attention dans
une époque particulièrement avare de ces expériences.

\bookmarksetup{startatroot}

\chapter{\texorpdfstring{\textbf{Chapitre 1. Pourquoi le geste
?}}{Chapitre 1. Pourquoi le geste ?}}\label{chapitre-1.-pourquoi-le-geste}

Un homme est couché dans un ruisseau, recroquevillé. Sa tête est en
contact avec la berge, déposée sur une pierre. Il a les yeux entrouverts
sans fixité. Il flotte ou repose sur le fond. Il est suspendu. La photo
est en noir et blanc, je la redécouvre.

Min Tanaka a passé de nombreuses années à s'entraîner dans les ruisseaux
en se laissant porter, former et déformer par les eaux glacées du Japon.
Il s'agit bien, pour Min Tanaka, d'un entraînement et d'un
apprentissage. Passer de longues heures à écouter l'écoulement de l'eau
et l'accompagner dans un dessaisissement, non pas du corps, mais du
mouvement volontaire. Ce n'est pas un abandon. C'est une écoute ouverte,
un dessaisissement actif où se tisse précisément une danse. Pour le
spectateur : une quasi-immobilité, un corps ballotté, abandonné au
courant. Pour Min Tanaka : un apprentissage en cours.

Cette inversion du regard est le premier geste de ce livre. Ce que nous
appelons habituellement mouvement, le déplacement observable, mesurable,
reproductible n'est que la partie émergée. Il y a autre chose qui
colore, oriente, donne sens au mouvement. Min Tanaka danse sous le
mouvement visible. Il soustrait pour laisser paraître un geste. Il offre
une minuscule chance au sens de devenir visible.

La question que pose cette image est aussi celle que pose toute séance
de Technique Alexander : qu'est-ce que nous cherchons, exactement, quand
nous pratiquons ? Nous pouvons venir avec des réponses très précises :
une douleur chronique, une rigidité, une qualité de présence entrevue un
jour. Nous venons souvent avec des convictions sur le corps et le
mouvement. L'apprentissage va souvent consister à défaire un certain
nombre de ces idées reçues, non pour les remplacer par d'autres, mais
pour offrir une expérience différente. Et cette expérience passe en
premier lieu par une compréhension différente de ce qu'est un geste.

C'est pourquoi ce livre commence par là. Pas par le corps, pas par la
posture, ni par l'habitude, l'inhibition, l'appréciation sensorielle ou
le contrôle primaire\ldots{} mais par le geste.

Le geste traverse notre langage avec une trompeuse familiarité. Nous
parlons de faire un geste, d'un beau geste, d'esquisser un geste. Cette
familiarité masque une ambiguïté fondamentale : le geste ne recouvre
jamais le seul mouvement du corps. Lorsque je saisis un verre d'eau,
ai-je effectué un mouvement ou accompli un geste ? La distinction peut
sembler bien fine et comme toute chose fragile, accessoire, mais elle
ouvre une brèche conceptuelle où se joue toute la singularité de ce que
la Technique Alexander explore, et que nous cherchons à déplier.

Hubert Godard propose cette distinction éclairante ; le mouvement relève
du déplacement biomécanique observable, trajectoire d'un membre dans
l'espace, activation de groupes musculaires, cinématique articulaire.
Une machine peut produire un mouvement : le bras d'un robot industriel
qui saisit une pièce sur une chaîne de montage effectue un mouvement
parfaitement reproductible, mesurable, analysable en termes de forces et
de trajectoires. Le geste, lui, appartient à un autre ordre de réalité.
Hubert Godard le définit comme ce qui s'inscrit dans l'écart entre le
mouvement visible et la toile de fond tonique du sujet, ce qu'il nomme
le pré-mouvement. C'est dans cet écart, dans cette coloration, que
réside l'expressivité du geste humain, dont est démunie la machine. Nous
ne faisons jamais deux fois le même geste et cette variabilité, est
paradoxalement ce qui en permet la lecture.~

Saisir un verre d'eau, en tant que mouvement, c'est une coordination
main-bras-épaule, une préhension palmaire, un ajustement de la force de
serrage, du rapport graivitaire. Trinquer avec ce même verre lors d'une
soirée d'anniversaire, c'est accomplir un geste. Trinquer, c'est faire
signe à l'autre, s'inscrire dans un rituel de convivialité, manifester
une intention sociale, habiter par des gestes, un monde rythmé et tissé
de liens. Le mouvement biomécanique est identique ou presque, lever le
verre, le déplacer vers celui de l'autre. Mais le geste configure un
monde : un monde où existe la célébration, la reconnaissance mutuelle,
le lien.

Le geste ne se contente pas d'agir sur un objet pré-donné dans un espace
neutre. Il fait paraître un monde. Cette formulation, empruntée à Emma
Bigé, n'est pas métaphorique. Elle désigne quelque chose de précis sur
le plan phénoménologique : avant que je me mette à écouter quelqu'un, il
n'y a pas encore de parole, il y a des sons. C'est le geste d'écouter
qui fait advenir la parole comme parole, qui constitue le silence comme
porteur de sens, qui ouvre un monde où quelqu'un dit quelque chose à
quelqu'un d'autre. De la même manière, le geste devient geste quand il
est porté et reconnu comme adresse par un autre ( présent ou absent ).
C'est cette même présence-absence qui organise et rythme la façon dont
j'écris sur mon clavier et qui contribue à organiser ma pensée comme
adresse.~

Pointer du doigt est peut-être l'exemple le plus saisissant. En tant que
mouvement, c'est une extension de l'index, une stabilisation du poignet,
une rotation de l'épaule. Mais pointer est un geste : il fait émerger un
objet dans le champ perceptif partagé, il dirige l'attention de l'autre,
il instaure une référence commune. Ce simple geste ressaisit l'oiseau
qui fraie dans le ciel et l'ami à mes côtés, il reconfigure après sa
réalisation, la scène en son entier. Il est comme le marqueur du verbe
en allemand, placé en fin de phrase. Ce n'est qu'une fois le geste
accompli, que le sens de toute la séquence se révèle rétroactivement,
que la relation entre les protagonistes et le monde se trouve nouée dans
et par le geste.~

Emma Bigé formule cela de manière saisissante : le geste nomme
l'indissociabilité de l'agir, de l'être et du sentir. Cette phrase
désigne une unité fondamentale que notre langage habituel tend à
fragmenter. Lorsque je touche quelque chose, je ne fais pas qu'agir sur
un objet, éprouver une sensation tactile, être présent à moi-même dans
cet acte, ces trois dimensions sont rigoureusement indissociables. Le
geste est précisément le nom de cette indissociabilité.

C'est ce qui rend la distinction mouvement/geste non pas théorique, mais
pratiquement décisive pour quiconque travaille avec la Technique
Alexander. Lorsque le professeur pose ses mains sur l'élève, il ne
touche pas des vertèbres cervicales ni des muscles. Il touche une
personne --- c'est-à-dire qu'il entre dans un monde relationnel où le
toucher est proposition, invitation, écoute. Ce toucher-là et le toucher
d'un mécanicien qui manipule des pièces ne mobilisent pas les mêmes
récepteurs tactiles, la même organisation gravitaire, ne déploie pas le
même imaginaire. Ils font littéralement advenir des mondes différents et
par là-même des usages différents de ces mondes.

Mais pourquoi le geste est-il le lieu de cette indissociabilité ?
Pourquoi pas le mouvement, l'action, l'acte ? La réponse est dans ce que
le geste porte d'archéologie et c'est ici que l'enquête doit aller plus
avant. Car, avant d'apprendre quoi que ce soit, avant même que quelque
chose comme un ``je'' ne se constitue, il y avait une indiscernabilité.
Nous en éprouvons encore le malaise parfois, quand assis confortablement
dans un train, le train voisin nous emporte tout entier dans son
sillage, nous dépossédant un instant de notre sol et de notre cohérence
spatiale chèrement acquise.~

Un nourrisson couché sur le dos, agite ses jambes : voit-il ses jambes
bouger, ou voit-il des «~objets-jambes~» qui se déplacent dans son champ
visuel ? Sent-il qu'il fait bouger ses jambes, ou sent-il que ses jambes
bougent, comme le mobile au-dessus de lui ? Il y a bien un « je », mais
pas au sens où l'adulte le vit -- un je embryonnaire, un je naissant,
mais sans la prétention d'un agent souverain. Il ne distingue pas encore
la source interne de la source externe du mouvement, l'agir de l'être
agi, le moi du monde. Je bouge. Je suis bougé. Le monde bouge. Ces trois
formulations décrivent le même événement, indémêlable.

Ce n'est pas une phase révolue, un stade primitif qu'on aurait dépassé.
C'est une couche qui reste là, sous les habitudes constituées, sous le
contrôle appris, sous la représentation de soi comme agent séparé du
monde qu'il habite. Et ce que les praticiens somatiques cherchent,
souvent sans le formuler ainsi, c'est précisément l'accès à cette
couche. Non pas pour y régresser, mais pour la retrouver comme
ressource. Parce que c'est dans ce nouage qu'un geste nouveau à la
chance de pouvoir émerger.~

Dans ce flux originel, dans cet entrelacs que Merleau-Ponty nomme la
chair, un seul invariant traverse absolument tout : la gravité . C'est
cette dernière qui fait de nous des être humains-humus, venant de la
terre comme berceau et fond indissociable de notre attachement.~ Qu'on
bouge ou soit bougé, qu'on agisse ou subisse, elle est là. Elle
organise. Elle ne peut être attribuée ni au dedans ni au dehors, elle
est toujours dans le pli. C'est cette constance absolue qui en fait
l'organisateur primordial de la différenciation. Comme le fond blanc
d'une page permet de voir les lettres noires sans jamais être vu
lui-même, la gravité permet de distinguer ce qui varie --- les
mouvements, les objets, les autres --- sans jamais être sentie
isolément, toujours à travers des variations toniques, posturales,
relationnelles.

C'est autour de cet axe que quelque chose comme un ``je'' commence à
s'extraire du flux. Non pas comme une substance qui se découvrirait
préformée, mais comme un pattern stable qui émerge progressivement des
coordinations répétées, coordinations où gravitaire, affectif, perceptif
et relationnel ne sont jamais trois choses séparées qui s'additionnent,
mais trois aspects d'un seul événement de la chair. Se redresser pour
voir quelqu'un, c'est simultanément négocier avec la gravité, exprimer
un accordage affectif, s'ouvrir au geste. Le ``je'' émerge dans ce
triple mouvement, pas avant lui.

Ce que ce développement éclaire rétroactivement, c'est la nature de ce
qu'on cherche en venant à une séance. On ne cherche pas à apprendre
comment bouger correctement. On cherche à retrouver accès à ce niveau où
organisme et monde ne sont pas encore antagonistes --- où le sol n'est
pas un obstacle mais un appui, où la gravité n'est pas une contrainte
mais une relation, où le mouvement n'est pas un effet de la volonté mais
une émergence de la chair.

Et la raison pour laquelle ce livre commence par le geste --- plutôt que
par le corps, la posture, la biomécanique --- c'est que le geste est
précisément le lieu où cette couche originaire reste accessible,
palpitante. Le mouvement peut se décrire sans sujet : cinématique,
forces, trajectoires. Mais le geste, jamais. Le geste est toujours
transitif. Il va toujours vers quelque chose, vers quelqu'un, vers un
monde. Il configure une relation. Et c'est dans la relation --- entre
soi et le sol, entre soi et l'autre, entre soi et la gravité --- que la
transformation est possible.

Min Tanaka immergé dans le ruisseau, explore la couche où les trois
formulations restent indémêlées : je bouge, je suis bougé, l'eau me
bouge. Et dans cet espace, quelque chose se danse parce qu'il a cessé de
séparer ce qui ne l'était pas au départ.

La Technique Alexander, vue depuis cette archéologie, n'est pas une
technique pour améliorer la posture. C'est une pratique qui apprend à
retrouver, geste par geste, cet espace de l'indiscernabilité créatrice.
Non pas en défaisant le ``je~», on ne peut pas et ce n'est pas le but,
mais en le rendant moins rigide, moins contrôleur, plus poreux à ce qui
émerge quand on cesse de «~tout diriger~».~

Les chapitres qui suivent déploient chacun un geste fondamental de cette
pratique. Non pas des techniques à maîtriser, mais des enquêtes à mener.
Enquêtes sur la façon dont toucher, peser, regarder, inhiber et
s'orienter peuvent rouvrir, chaque fois, l'espace entre le mouvement
visible et la toile de fond tonique qui le porte. L'espace où, pendant
un instant, on ne sait plus tout à fait qui bouge et où quelque chose
peut commencer, re-commencer.

\bookmarksetup{startatroot}

\chapter*{References}\label{references}
\addcontentsline{toc}{chapter}{References}

\markboth{References}{References}

\protect\phantomsection\label{refs}




\end{document}
